\documentclass[11pt,a4paper,twocolumn]{article}
\usepackage[T1]{fontenc}
\usepackage{mathptmx}
\usepackage[left=18mm,right=18mm,top=18mm,bottom=19mm,columnsep=6mm]{geometry}
\usepackage{amsmath,amssymb}
\usepackage{algorithm,algpseudocode}
\usepackage{microtype,titlesec}
\usepackage[nospread]{flushend}
\usepackage[hidelinks]{hyperref}
\titleformat{\section}{\large\bfseries}{\thesection}{0.6em}{}
\titleformat{\subsection}{\normalsize\bfseries}{\thesubsection}{0.6em}{}
\titlespacing*{\section}{0pt}{8pt plus 2pt minus 1pt}{4pt}
\titlespacing*{\subsection}{0pt}{6pt plus 2pt minus 1pt}{3pt}
\setlength{\parindent}{1em}
\setlength{\parskip}{1pt}
\setlength{\textfloatsep}{8pt plus 2pt minus 2pt}
\setlength{\intextsep}{8pt plus 2pt minus 2pt}
\setlength{\abovedisplayskip}{5pt plus 2pt minus 1pt}
\setlength{\belowdisplayskip}{5pt plus 2pt minus 1pt}
\renewcommand{\topfraction}{.92}
\renewcommand{\bottomfraction}{.85}
\renewcommand{\textfraction}{.08}
\setcounter{topnumber}{3}
\emergencystretch=1em
\algrenewcommand\algorithmicrequire{\textbf{Input:}}
\newcommand{\para}[1]{\par\noindent\textbf{#1}\enspace}
\newcommand{\TTFT}{\mathrm{TTFT}}
\newcommand{\TPOT}{\mathrm{TPOT}}
\hypersetup{pdftitle={PDblend: Profiling, Online Scheduling, and Adaptive Pool Configuration}}

\begin{document}
\pagestyle{plain}
\raggedbottom
\twocolumn[{\centering\Large\bfseries
PDblend: Profiling, Online Scheduling,\\
and Adaptive Pool Configuration\par\vspace{9pt}}]

PDblend combines mixed execution (M) with prefill/decode
disaggregation (P/D) within a GPU cluster. It seeks to reduce
serving energy while satisfying time-to-first-token (TTFT) and
average time-per-output-token (TPOT) objectives. An
\emph{instance} comprises an inference engine and its tensor
parallel (TP) GPU group. An M instance executes both stages;
a PD request uses a pair of P and D instances. The method has
three components: an offline profiler provides cost estimates,
an online router assigns requests to execution paths, and an
adaptive controller adjusts frequencies, roles, and traffic
shares across resident TP pools. The following description
reflects the current implementation and identifies boundaries
where its execution paths differ.

\section{Profiler}
\label{sec:profiler}

The profiler estimates the cost and capacity of an instance
under a given model, TP layout, frequency, and request shape.
Each topology uses an independent profile because changing TP
affects communication, model memory, and KV capacity.
Comparisons across frequencies likewise depend on the
supported measurement domain. Model, engine, hardware, and
measurement identities bind a profile to its intended
deployment.

\para{Measurement dimensions.}
Prefill measurements vary the prompt length $L$ and GPU
frequency $f$. Decode measurements additionally vary batch
size $B$ and context length $C$, recording the shapes actually
executed during generation. Mixed measurements observe
prefill alongside active decodes, capturing the stalls that
long prompts impose on ongoing requests. Collection separates
loading, warmup, and clock setup from the measurement windows,
and preserves raw observations and independent validation
samples. Runtime components also cover idle, parked, and
stopped states, wakeup and clock-change latency, KV capacity,
and P/D handoff.

\para{Time and power models.}
Let $T_p(L,f)$ denote prefill time and $T_d(B,C,f)$ denote
decode step time. The current native timing model uses a
nonnegative fit at each measured frequency:
\begin{align}
T_p(L,f)&=a_f+b_fL+c_fL^2,\\
T_d(B,C,f)&=u_f+v_fB+w_fBC.
\end{align}
The linear terms represent the main workload contribution,
while $BC$ captures context-dependent memory access.
Normalization used in the implementation is absorbed into
the coefficients. Compatibility profiles may also use a
$B^2$ term or piecewise curves; the bound profile version
determines the actual query behavior. Power models depend on
frequency and batch/context shape. Where supported by
measurements, mixed power comes from a mixed serving window.

Timing and power must refer to consistent observation
boundaries. For a serving window of duration $H$,
\begin{equation}
E_{\mathrm{window}}=\int_0^H P(t)\,\mathrm{d}t=H\bar P,
\end{equation}
where $\bar P$ is its mean power. A native request-cycle mean
includes scheduling and intervals between kernels. Multiplying
that mean by CUDA-only prefill time would mix incompatible
boundaries. Pool energy therefore uses a corresponding cycle
or layout-window model. Parked capacity still contributes
static power.

\para{Transfer and validated queries.}
The basic KV cost interface uses a fixed term plus data volume
divided by effective bandwidth:
$T_x(L)=x_0+Lk/v$, where $k$ is KV bytes per token and $v$
is effective bandwidth. Other native components interpolate
observed values at measured prompt lengths. A difference
between M and PD path latency, however, includes protocol and
scheduling effects and is not a direct physical-copy
measurement. Section~\ref{sec:online} instead defines handoff
through observable output-token endpoints.

Calibration and tuning data support fitting and selection;
independent holdouts check prediction error and selected
configurations. Published profiles carry frequency, length,
batch/context, and topology domains. Bounded versions reject
queries outside their supported coverage. Interpolation does
not by itself establish validation for an additional domain.
Before planning, the consumer also checks that the required
capacity, static-state, transfer, and clock-transition
components are available. A profile thus supplies a prediction
together with its identity, measurement boundary, and range
of applicability. Decisions retain the profile binding for
subsequent inspection.

\section{Online Algorithm}
\label{sec:online}

For each arriving request, the router fixes its prefill and
decode owners. Let $L$ be its input length and $K$ its maximum
output length. An M route assigns both stages to one instance.
A PD route selects an ordered pair $(p,d)$: $p$ processes the
prompt and $d$ continues generation. Both pools may contain
multiple instances; the pair is selected for the request,
rather than fixed for the entire cluster.

Instance-local token batching remains the engine scheduler's
responsibility. The default waiting policy is FCFS with
continuous batching and chunked prefill, so admission order
does not imply completion order. Within a TP instance, GPUs
jointly execute the scheduled batch; they do not independently
accept requests as individual GPUs become idle.

\subsection{Candidate Paths and Instance Pairing}

The router reads the published role and admission state of
each instance. A PD pair must contain distinct instances with
matching model, TP, pipeline parallelism (PP), pool identifier,
and generation. The current path supports PP=1 and does not
remap KV shapes across different TP layouts. The router
therefore enumerates compatible pairs before choosing their
members.

When a compatible pair exists, a prompt with $L\ge\tau$
normally prefers PD; other prompts prefer M. If no M instance
accepts requests, a short prompt may also use PD. Without a
complete PD pair, a long prompt can use M. An enabled pressure
policy imposes its own minimum PD input length as well.
The threshold $\tau$ selects a preferred branch, rather than
unconditionally determining the final route.

For instance $i$, the proxy maintains $U_i$, the total prompt
tokens assigned to requests whose first token has not yet
been observed, and $N_i$, the number of assigned, unfinished
decode sequences. Dispatch immediately increases $U_p$ and
$N_d$, allowing later decisions to see future decode load
even while prefill is running. These are proxy estimates:
$U_i$ is not an exact per-chunk remainder, and $N_i$ is not
the engine's running batch size. The basic load rule compares
$(N_i,U_i)$ lexicographically for M and $(U_p,N_d,U_d)$ over
compatible PD pairs.

The standard runner enables an additional SLO routing layer.
It first checks request length, admission state, and KV
capacity for all candidates in the preferred branch.
Each endpoint of a PD route is charged the $L+K$ request
budget in proxy admission accounting; the engine allocates
the actual KV blocks. The router then queries the profile at
the current control frequencies. Predicted P queueing sums
the prefill costs of waiting requests individually. Decode
cost depends on assigned sequences and their context bounds.
Summation avoids treating several queued prompts as one
fictitious prompt outside the profiling domain.

\subsection{Charging the KV Handoff}

The proxy first submits the PD request to P. After prefill,
P generates a first token, which the proxy forwards to the
client. The proxy then submits continuation to D with the
carried token, transfer identity, and KV metadata, reducing
the remaining output budget to $K-1$. The connector may
submit KV layer by layer during P execution, so physical
transfer can start before P emits its token. D continues
once the required KV is available. A one-token PD candidate
uses only P and does not initiate remote KV transfer.

\begin{samepage}
Let $Q_p$ be P queueing time,
$s=T_d(B,C,f_D)$ the D step time, and $h$ the interval
from P's first token to D's first token. For $K\ge2$, the
ordinary router predicts
\begin{align}
\TTFT_{PD}&=Q_p+T_p(L,f_P),\\
\TPOT_{PD}&=\frac{h+(K-2)s}{K-1}.
\label{eq:handoff}
\end{align}
\end{samepage}
The interval $h$ includes remaining KV handoff, proxy
submission, D queueing and reception, and the first D step.
An observed endpoint interval already includes that step;
adding $s$ again would double-count it. When $K=2$, the
entire TPOT budget is consumed by this interval. Longer
outputs amortize the one-time handoff over subsequent tokens.
Prompt length alone is consequently insufficient to judge
whether PD is appropriate.

For M, the basic estimate is
$Q_m+T_p(L,f_M)+T_d(B,C,f_M)$. The router additionally
checks whether an incoming prefill would violate the
first-token or remaining output budgets of incumbent
requests. This prevents spillover from protecting a new
request at the expense of active decodes.

For outputs longer than two tokens, missing endpoint
coverage falls back to
$h=\max(T_x(L)+s,h_{\min})$.
The configured floor $h_{\min}$ is a risk parameter, not
measurement evidence. A two-token output requires matching
endpoint coverage before PD can be classified as SLO-safe.
The current development endpoint component covers batch 1,
16 output tokens, requested P/D clocks of 2520~MHz, and
inputs of 512--7168 tokens. It interpolates training maxima;
these observations do not validate two-token replies,
loaded decode batches, or tail guarantees. Online timestamps
record actual handoff and progress without automatically
extending this offline coverage.

\begin{algorithm}[t]
\caption{Request assignment in the ordinary router}
\label{alg:route}
\small
\begin{algorithmic}[1]
\Require Request $(L,K)$, roles, profile, and SLO
\State Build preferred candidates using $\tau$ and compatibility
\State If PD is preferred, also enumerate M alternatives
\State Remove candidates failing length, admission, or KV checks
\For{each remaining candidate}
  \State Predict queueing, compute, handoff, and incumbent risk
  \State Mark profile coverage and TTFT/TPOT feasibility
\EndFor
\If{the preferred branch has an SLO-safe candidate}
  \State Choose the candidate with lowest predicted TTFT
\ElsIf{PD is preferred and an SLO-safe M exists}
  \State Choose the safe M with lowest predicted TTFT
\Else
  \State Retain an admitted original route; mark SLO unproven
  \State Return no route if the original branch has no capacity
\EndIf
\State Pin $(p,d)$, update $U_p,N_d$, and start execution
\end{algorithmic}
\end{algorithm}

\subsection{Fallback and Cost-Model Boundaries}

Algorithm~\ref{alg:route} selects the lowest predicted TTFT
among preferred candidates meeting both latency budgets with
a safety margin. It spills PD to M only when the preferred
PD branch has no safe candidate and a safe M alternative
exists. If predictions are unavailable or unsafe, the current
implementation may retain a capacity-admitted original path,
explicitly marking its SLO status as unproven. Admission is
therefore not an unconditional SLO guarantee. Dispatched
requests retain their owners. For an abnormal termination
after engine submission, uncertain cleanup must be resolved
by native acknowledgement before accounting is released.

The endpoint formulation above applies to the ordinary
router. Pool-level planning still adds the legacy transfer
estimate to PD TTFT. The default ResidentRouter score has no
explicit transfer term; its optional SLO/energy context adds
transfer to first-step latency and has not adopted
Equation~\eqref{eq:handoff}. With an explicitly configured,
validated incremental-energy estimator, resident routing
can compare added joules after coverage, latency, and
capacity checks. If any retained alternative lacks energy
evidence, the entire set falls back to time scoring.
There is no separately validated pure KV-copy energy term;
matched serving windows or measured whole-route increments
account for its aggregate effect.

\section{Adaptive Pool Configuration}
\label{sec:adaptive}

While routing handles individual arrivals, adaptation
determines capacity over the next planning horizon.
The controller uses arrival rates, input/output lengths,
and outstanding work to choose P/D/M counts, role
frequencies, parked states, and the PD/M threshold.
Writing a configuration as $c$ and the current configuration
as $c_0$, the objective is
\begin{equation}
c^*=\arg\min_{c\ \mathrm{feasible}}
       \{H P(c)+E(c_0,c)\}.
\label{eq:pool-objective}
\end{equation}
Here $P(c)$ includes active and parked power, while
$E(c_0,c)$ is measured or estimated transition energy.
Feasibility requires latency margins, KV capacity, profile
coverage, and preservation of branches owning outstanding
requests. Each candidate is charged its own transition
cost. If the current configuration remains feasible,
a change must exceed a relative saving threshold.
The default amortization horizon is 60 seconds; the standard
PDblend policy sets the saving threshold to 8\%.
Runtime configuration may override both.

Transition costs preferably come from paired measurements
bound to the source/target configurations and profile.
A strict mode excludes uncovered transitions; the ordinary
fallback estimates wakeup and clock-change costs without
labeling them as measurements. Forecasts retain waiting
prefill tokens, remaining decode work, and occupied KV,
including long requests that have left the recent statistics
window. The controller normally observes protection signals
every second and replans every ten seconds. Missing coverage
leads to retention or conservative recovery rather than
qualifying an unsupported energy optimization.

\subsection{Frequency Control}

Frequency control changes service rate while preserving the
model and TP topology. The planner enumerates supported
discrete frequencies and selects $f_P,f_D,f_M$ using
stage-specific load, latency, and power estimates.
P and D may use different clocks, whereas all GPUs within
one TP instance receive the same target frequency.
The search is discrete enumeration; physical actions only
update instances whose state changes.

When the budget-aware Shield is enabled, its fast path
identifies the pressured role: PD prefill pressure targets P,
decode pressure targets D, and mixed-path pressure targets M.
It first raises relevant clocks, then restores capacity
under persistent pressure. A brief collective token gap
permits a bounded clock probe rather than immediately
establishing a capacity shortage. Optional deadline
protection searches higher frequencies for the lowest
one that removes the predicted risk and maintains an M
frequency floor until repeated safe observations permit
release. Hold times, cooldowns, and repeated downshift
confirmation suppress oscillation.

Hardware actions use physical-GPU locks so concurrent
controllers cannot write the same device simultaneously.
Blocking operations run in worker threads, allowing urgent
clock actions on other instances while one engine loads.

\subsection{Role Change}

Role control redistributes instances among P, D, and M,
subject to the fleet budget and service constraints.
Unused instances can remain resident in a low-power state
or have their processes stopped. When mapping target counts
to concrete instances, the controller first preserves
existing roles where possible, then considers residency
and load to avoid unnecessary wakeups and reassignment.

An active-to-active role change updates the role and
frequency used for subsequent arrivals. Existing requests
remain pinned to their original instances, with no migration
of live KV. The new role therefore takes effect gradually:
an instance can still be finishing work admitted under its
previous role. Forecasts retain the original branch
ownership of that backlog.

Parking or stopping has a stronger boundary. The controller
first disables new admission, waits for proxy-side prefill
and decode work to drain, and, when native control is
connected, verifies every rank's queues, KV blocks, and
pending transfers before reducing power or stopping.
Wakeup performs startup or resume, readiness checks,
clock configuration, and native admission restoration
before publishing availability. Actions on different
instances can proceed concurrently, but failures prevent
affected instances from accepting requests.

The standard policy retains at least $\min(4,N)$ M
instances, where $N$ is the pool's instance count.
Its default search consequently does not admit arbitrary
pure-PD layouts. Lower floors require acceptance evidence
matching the model, topology, SLO, and workload domain;
specific runners additionally constrain frequencies and
reserve recoverable M capacity. This guards against
over-consolidation when the queueing approximation cannot
fully capture mixed interference.

\begin{algorithm}[t]
\caption{Adaptive pool control}
\label{alg:adapt}
\small
\begin{algorithmic}[1]
\Require Current configuration, profiles, history, and backlog
\State Update forecasts and backlog, preserving owners
\State Observe latency and output progress in the fast loop
\If{latency risk is detected}
  \State Raise clocks; restore capacity for persistent pressure
  \State Enforce the protection hold period
\EndIf
\If{replanning is triggered and optimization is allowed}
  \State Enumerate roles, clocks, parked states, and thresholds
  \State Enumerate outer TP shares if joint mode is enabled
  \State Filter candidates by coverage, capacity, and SLO
  \State Minimize horizon energy plus transition cost
  \State Apply saving and stability gates
  \State Execute changes while preserving request owners
  \State Publish TP shares only after all pools are ready
\EndIf
\State Record actual state and dispatch counts for feedback
\end{algorithmic}
\end{algorithm}

\subsection{TP Weight Change}

TP weights denote request shares across two resident pools
with different TP sizes. The pools occupy disjoint GPUs
and use independent profiles. Online changes preserve
their existing model-shard layouts. An outer search
enumerates shares $w$ and $1-w$; each pool's inner planner
chooses roles and frequencies. Their energies are combined
using Equation~\eqref{eq:pool-objective}. The default share
grid has step 0.1 and also includes the current target
and the preceding interval's observed allocation.

Future arrival rates are scaled by the candidate shares,
but in-flight work remains charged to its original pool.
Even a pool assigned zero new traffic retains its full
resident-capacity cost, and unused GPUs remain in power
accounting. Both inner plans must satisfy coverage,
capacity, and SLO constraints. If the current joint
configuration is still feasible, the proposal must also
meet the saving threshold. New shares are published only
after all pools' roles, frequencies, and admission states
match the completed plans.

Joint planning requires measured decode and mixed-window
power. Shield escalation or a protection hold suspends
energy optimization. If any pool fails during execution,
the coordinator retains the old target shares and blocks
further optimization pending state recovery.
Algorithm~\ref{alg:adapt} summarizes the shared control
structure; it does not prescribe a fixed sequence in which
all three kinds of change must occur.

Dispatch realizes the target through weighted deficits.
After $A$ assignments, let pool $j$ have received $a_j$
requests. Among currently feasible pools, it selects
\begin{equation}
j^*=\arg\max_j\{w_j(A+1)-a_j\}.
\end{equation}
If the target pool is infeasible, another feasible pool
can serve the request. Shares are thus allocation targets,
not mandatory rejection rules. Actual dispatch counts feed
the next planning interval, exposing deviations caused by
capacity constraints. This mechanism does not migrate
in-flight requests across TP layouts. Online model-weight
resharding remains disabled because a GPU-validated native
transaction backend is not yet available.

\end{document}
